\documentclass[11pt]{amsart}

\usepackage{amsmath, amssymb}
\usepackage[margin=1in]{geometry}
\usepackage{enumitem}
\usepackage{xcolor}
\usepackage[colorlinks=true, linkcolor=blue, urlcolor=blue]{hyperref}

% ---------------------------------------------------------------------------
% Template macros.
%
%   \todo               placeholder for an unwritten response
%   \response{...}      the response block attached to a referee request
%   \refitem{tag}{...}  a referee request, tagged so it can be \cref'd/\ref'd
%
% To fill in the file, replace each \response{\todo} with \response{<text>}.
% Grep for "TO DO" to find what is left.
% ---------------------------------------------------------------------------

\newcommand{\todo}{\textcolor{red}{\textbf{TO DO}}}
\newcommand{\done}{\textcolor{green}{\textbf{DONE}}}

\newcommand{\response}[1]{%
  \par\smallskip
  \noindent\hangindent=2em\hangafter=1
  \textit{Response.}\ #1
  \par\smallskip}

% Referee items are labelled R1.n / R2.n / R2.Tn so that the numbering can be
% quoted back to the referees and cross-referenced between the two reports.
\newlist{requests}{enumerate}{1}
\setlist[requests]{leftmargin=3.2em, itemsep=0.8em, label=\textbf{\rlabel.\arabic*}}
\newcommand{\rlabel}{R}

\title{Response to referees}
\author{}
\date{\today}

\begin{document}

\maketitle

\noindent
We thank both referees for their careful and constructive reading of the paper.
Below we list each request in turn and describe the corresponding change to the
manuscript. Referee~1's items are numbered \textbf{R1.$n$} following the order of
that report; Referee~2's items retain that report's own numbering, as
\textbf{R2.$n$} for the numbered suggestions and \textbf{R2.T$n$} for the typos and
minor comments.

\medskip

\noindent
\textbf{Overlapping requests.} The following are raised by both referees and should
be answered consistently:
\begin{itemize}[leftmargin=2em, itemsep=0.3em]
  \item The status of the non-formalized extension of the generating set to the full
        implication graph, and whether the external transitivity/duality tools can be
        trusted: \textbf{R1.1}--\textbf{R1.6} and \textbf{R2.4}.
  \item Compilation timing data for the full versus the reduced set:
        \textbf{R1.7} and \textbf{R2.4}.
  \item The trusted axiom set / trusted computing base: \textbf{R1.3} and \textbf{R2.3}.
  \item Section~11 (the CNN section): \textbf{R1.18} (shorten and move to an appendix)
        and \textbf{R2.13} (justify the CNN heuristic against simply running an ATP,
        and strengthen the evidence against transitivity-learning).
  \item The phrase ``proof assistant language Lean'': \textbf{R1.M1} and \textbf{R2.T1}.
\end{itemize}

% ===========================================================================
\section*{Referee \#1}
\renewcommand{\rlabel}{R1}

Referee~1 recommends \emph{major revisions}, raising five major issues together with
three miscellaneous comments.

\subsection*{Exact theorem and trust base}

\begin{requests}

\item The abstract (``all validated by the formal proof assistant language Lean''),
the introduction (``the \(4694 \times 4693 = 22\,028\,942\) implications \dots\ were
completely determined, with proofs or refutations formalized in Lean'') and
\S1.3 Outcomes (``with all of the 22\,033\,636 implications \dots\ formalized in the
proof assistant language Lean'') all give the impression that the full implication
graph was formalized in Lean. The referee regards the \S1.3 sentence as simply wrong.
These passages must be corrected.
\response{\todo}

\item The decision to formalize only a reduced generating set (10\,657 positive and
586\,925 negative implications), with the extension to the remaining implications
performed by programs external to Lean, is currently explained only on page~12. The
referee considers the decision itself entirely sound, but insists that it be stated
early in the paper rather than deep in the technical sections.
\response{\todo}

\item Give a clear statement of exactly what must be trusted in order to believe that
all 22\,028\,942 implications are settled. The referee's reading is that the trusted
base consists of (i)~the Lean kernel, (ii)~the Lean program extending the generating
set to the full graph, and (iii)~some process confirming that all 22M results are in
fact settled; this should be confirmed or corrected explicitly.
\response{\todo}

\item Clarify the role and format of \texttt{full\_entries.json}. The referee expected
it to be the artifact underlying point~(iii) above, but found its
\texttt{"proven": false} entries confusing.
\response{\todo}

\item Provide instructions enabling an independent reader to verify the claim for
themselves --- specifically, how one would write a hand-written checker that deduces
the full implication table from the Lean-verified results. The referee states that in
the paper's current state they cannot gain full confidence that all 22\,028\,942
implications are settled.
\response{\todo}

\item Discuss which shortcomings of Lean, or of proof assistants in general, made the
split into a Lean-formalized generating set plus an unverified deduction of its
consequences necessary. (See also \textbf{R1.20}--\textbf{R1.23}.)
\response{\todo}

\end{requests}

\subsection*{Discussion of the Lean 4 artifact}

\begin{requests}[resume]

\item Report the extent of the final formalization, at minimum in lines of code and
compile times, and identify which parts dominate the compile time --- the
ATP-derived parts or the manual ones.
\response{\todo}

\item Explain how the millions of theorem statements were generated, and why that
generation process can be trusted.
\response{\todo}

\item State which sections of the paper correspond to formalized results and which do
not (the referee notes \S10 as an example of the latter).
\response{\todo}

\item Reference a specific tag or commit of the repository described in the paper,
and make clear which Lean and Mathlib versions are intended.
\response{\todo}

\end{requests}

\subsection*{Discussion of the collaboration}

\begin{requests}[resume]

\item The paper states as its goal ``to explore new ways to collaboratively work on
mathematical research projects using machine assistance'', but does not sufficiently
analyse those new ways. A detailed comparison with other projects is requested.
\response{\todo}

\item Discuss the Liquid Tensor Experiment. The referee notes that it used the same
organization as the PFR project, predates it, and was as far as they know the first
project centred around blueprints; it was likewise coordinated via Zulip and relied
on blueprints, but was more centralized and more clearly hierarchical than the ETP.
Without a discussion of the differences, the referee suggests the accurate framing is
that the ETP experimentally confirmed the model used in the LTE.
\response{\todo}

\item Treat the Busy Beaver Challenge (BB(5)) more accurately: the referee finds the
present reference inaccurate, since that project was less about the use of proof
assistants (it was developed with computer assistance but largely without proof
assistants, with formalization following after the fact). Discuss how its coordination
choices --- Discord and a wiki-like website --- compare with the ETP's, given that the
paper sets the exploration of new ways of collaborating as its prime target.
\response{\todo}

\item Discuss earlier large collaborations: Flyspeck and the odd order theorem, and in
a less ``named'' sense the AFP and math-comp. The referee also asks what differences
were found between a large mathematical collaboration and the large collaborations
typical of proof assistant and software development.
\response{\todo}

\item Provide more quantitative data on the collaboration itself: how the set of
contributors varied over the course of the project, the shape of the curve along which
results ``fell'', and how the two are correlated.
\response{\todo}

\end{requests}

\subsection*{Structure of the paper}

\begin{requests}[resume]

\item The referee finds the structure chaotic, reading as separately written sections
with no clear order, and the table of contents as mixing mathematics, methods and
tools. Reorganize into a clear part on the mathematics, a second on the methods of
collaboration, and a third on the tools used.
\response{\todo}

\item Section~7 is 13 pages long, with \S7.3 alone accounting for 9 of them. Move the
most technical material --- \S\S7.3.2, 7.3.3, 7.3.5, 7.3.6, 7.3.9 and Example~7.1 ---
to an appendix.
\response{\todo}

\item Section~11 contains no formalized results and reports mainly preliminary work,
since the randomized control experiment was not run. Shorten it and move it to an
appendix. (See also \textbf{R2.13}.)
\response{\todo}

\item The observation that using strong automation early in a project kills insights
is, in the referee's view, one of the most interesting remarks in the paper, but is
buried in \S14.1. Give insights of this kind greater prominence.
\response{\todo}

\end{requests}

\subsection*{Discussion of the choice of Lean 4}

\begin{requests}[resume]

\item Address the question from \emph{Virtues of a Formalization Project}: did the
choice of proof assistant or logical foundation crucially influence the formalization
process? The referee regards this as expected of an AfM paper, and identifies three
factors: ATP capabilities, native computation, and mathematical libraries.
\response{\todo}

\item On ATP capabilities: compare with Isabelle, where \texttt{sledgehammer} calls
provers such as CVC4 and Vampire natively and proof reconstruction is natively
supported, against the custom reconstruction architecture that had to be written here.
\response{\todo}

\item On native computation: the paper notes that \texttt{native\_decide} was avoided
because of its known soundness issues, but the referee would like to see how much
\texttt{native\_decide} would have sped up the counterexample computations of
Figures~5 and~6 relative to \texttt{decide} and \texttt{decideFin!}, ideally with a
build-time comparison.
\response{\todo}

\item Consider whether the external reconstruction of the 22M implication matrix could
instead be validated inside a proof assistant using native computation, and in
particular whether verifying the matrix in Rocq would be feasible.
\response{\todo}

\item On mathematical libraries: address the referee's impression that the Mathlib
aspect --- usually the key reason Lean~4 is chosen --- did not actually matter for
this project.
\response{\todo}

\end{requests}

\subsection*{Miscellaneous comments}

\begin{requests}[label=\textbf{R1.M\arabic*}]

\item (p.~4) ``proof assistant language Lean'' is awkward; the referee suggests
``in the language of the proof assistant Lean''. (See also \textbf{R2.T1}.)
\response{\todo}

\item (p.~11) The explanation of the custom commands (\texttt{equation},
\texttt{EquationX}, \texttt{LawX}) cannot be followed; examples would help.
\response{\todo}

\item (p.~18) The paragraph beginning ``ideally the environment'' needs more
elaboration; the referee does not understand the point it is making, and in particular
whether ``practical implementations'' refers to more than just Lean.
\response{\todo}

\end{requests}

% ===========================================================================
\section*{Referee \#2}
\renewcommand{\rlabel}{R2}

Referee~2 recommends \emph{minor revisions} and gives the paper their highest
recommendation for acceptance, subject to the comments below.

\subsection*{Suggestions}

\begin{requests}

\item Cite tools at their first mention, not only where they are discussed topically.
Lean is first mentioned on p.~2, line~9 but first cited on p.~17; the same applies to
the Lean blueprint tool (first appearing p.~12, line~11) and to \texttt{egg} and
\texttt{duper} (p.~10, line~26).
\response{\todo}

\item Add citations for tools and libraries that currently lack them: Mathlib
(p.~10, line~24; see \url{https://leanprover-community.github.io/cite.html}),
SAGE (p.~45, line~36; see \url{https://wiki.sagemath.org/Publications_using_SageMath}),
GAP (p.~45, line~37; see \url{https://www.gap-system.org/cite/}), and the GAP small
group library (p.~45, line~37; see \url{https://gap-packages.github.io/smallgrp/#citing}).
\response{\todo}

\item (p.~11, lines~17--18 and p.~12, lines~1--2) The paper states that
\texttt{lean4checker} was used to limit the axioms used to a ``small trusted set''.
Say explicitly which axioms are in that set, and whether the consistency of that set
is well known. (See also \textbf{R1.3}.)
\response{\todo}

\item (p.~12, lines~3--9) Include timing data for compiling the full and the reduced
sets, so the reader can judge whether the difference is minutes, hours or days, and
say whether any effort has been made --- or any avenue investigated --- to speed the
compilation up. The referee further notes that the extension to the full graph via
transitivity and duality is done by external tools, and asks whether those tools are
correct and can be trusted; they find it strange that this is not treated with the same
care as the deliberate avoidance of \texttt{native\_decide} in the preceding paragraph.
(See also \textbf{R1.1}--\textbf{R1.6} and \textbf{R1.7}.)
\response{\todo}

\item Replace the unlocated references to the ETP blueprint with citations to specific
chapters or sections, since the blueprint is a standalone document of 27 chapters
across 128 pages. The affected places are footnote~4 on p.~3; p.~7, line~37; p.~8,
lines~11 and~20; p.~9, line~5; p.~30, line~20; p.~33, line~18; p.~52, lines~19 and~30;
and p.~67, line~13. The explicit URLs to blueprint chapters on p.~64, lines~15 and~40
should likewise be replaced by proper citations.
\response{\todo}

\item (p.~26, lines~6--7) The paper says the functional equations become too complex
for \(n > 2\). Clarify whether this exploration was carried out by hand or with ATPs or
other automation, and whether ``complex'' is meant informally (too complex for human
analysis) or formally (belonging to some intractable complexity class).
\response{\todo}

\item (p.~29, lines~5--7) The paper says the procedure for finding a rule set for the
greedy extension method can be automated using ATPs or SAT solvers. State explicitly
whether this automation was implemented as part of the project, and whether a Lean
tactic or external program for it exists.
\response{\todo}

\item (Remark~6.5, p.~35, line~7) The remark observes that matching invariants can come
from other algebras with a larger signature. State whether such matching invariants
were used in the formalization, and if so give an example.
\response{\todo}

\item (footnote~14, p.~43) The footnote notes that the timings in \S7.3 were performed
in heterogeneous environments and cannot be taken as benchmarks, but the hardware used
is never stated. State the hardware, and make explicit which experiments were run on
the same hardware. The same applies to the ``165 CPU-hours'' search reported on p.~21,
line~2.
\response{\todo}

\item (footnote~22, p.~53) Provide evidence for the claim that this is the
lowest-numbered law that does not have full spectrum but does have models of sizes
2 and~3.
\response{\todo}

\item (p.~54, lines~7--9) Cite the article in preparation in which the results on
spectra of laws of higher orders will be reported. The same applies to the claims about
modules on p.~54, lines~25--27, and to the extensions mentioned on p.~53, line~3.
\response{\todo}

\item (footnote~25, p.~57) Describe the format and meaning of the data in the linked
JSON file, which is otherwise not straightforward to interpret.
\response{\todo}

\item (\S11) Make explicit what advantage the CNN heuristic has over simply running
Prover9 or Vampire~5.0 with the previously determined optimal parameters for a fixed
time and assuming failure means non-implication --- an approach which, given the
paper's own statements on p.~48, lines~1--9, has 100\% accuracy up to order~4. The
referee also finds the evidence that the CNN is not learning some form of transitivity
inconclusive, and suggests either performing the additional tests described on p.~58,
lines~24--26, or analysing the CNN from an explainability and interpretability
standpoint. (See also \textbf{R1.18}.)
\response{\todo}

\item (p.~64, lines~6--7) The remark that ATPs outperformed LLMs ``if one restricts to
publicly available LLMs'' implies the existence of non-public LLMs that outperform
ATPs, for which no evidence is given. Either cite the publication establishing this, or,
if the qualification was made out of caution, list explicitly which LLM tools were tried
and did not yield results.
\response{\todo}

\end{requests}

\subsection*{Typos and minor comments}

\begin{requests}[label=\textbf{R2.T\arabic*}]

\item (Abstract, line~5) ``all validated by the formal proof assistant language Lean''
is confusing: validation is performed by the Lean kernel, whereas the Lean language is
that in which the proofs are written and performs no validation itself. The referee
suggests ``all validated by the formal proof assistant Lean'' or ``all validated by the
interactive theorem prover Lean'', to match the terminology of \S4.7. (See also
\textbf{R1.M1}.)
\response{\todo}

\item (p.~3, line~4) ``In this entailment pre-ordering'' \(\to\) ``In this entailment
pre-order''.
\response{\todo}

\item (p.~4, line~24) ``contributions had to be entered in the proof assistant language
Lean'': ``entered'' is clunky; ``written'' or ``formalized'' would be better.
\response{\todo}

\item (p.~4, line~30) ``written in the proof assistant language Coq'': since p.~17,
line~12 notes the name change (``Rocq (formerly Coq)''), make that note here at the
first occurrence and use ``Rocq'' thereafter.
\response{\todo}

\item (p.~5, line~10) ``either 820 or 822 pairs'' is confusing. If the discrepancy is
due to the two implications of unknown status, say so explicitly --- that the finite
implication held for 820 additional pairs and that for 2 pairs the status is unknown.
\response{\todo}

\item (p.~10, line~21) ``were also first generated as computer output from a source
other than Lean'': ``generated'' already implies ``as computer output''. Consider
``were also first generated from a source other than Lean''.
\response{\done}

\item (p.~10, line~37) ``containing two words in a free group'' should be ``in a free
magma'', per the definition of \texttt{EquationX} in the Lean sources.
\response{\todo}

\item (p.~12, line~11) ``LeanBlueprint'' is used, but p.~13, line~22 and thereafter use
``Lean blueprint''. Use one form consistently.
\response{\todo}

\item (p.~15, line~30) ``leanchecker'' is used, but p.~11, line~18 and p.~18, line~15
use ``lean4checker''. Clarify whether these are the same tool or a typo; if they are
genuinely different, cite each.
\response{\todo}

\item (p.~17, lines~9 and~34) ``relatively spartan'' and ``spartan language'': the
meaning is unclear, and if this is not technical jargon with a specific meaning another
adjective would be more appropriate.
\response{\todo}

\item (p.~33, line~15) ``versions of the greedy algorithm are useful'': if this refers
to the abstract greedy algorithm of Theorem~5.12, reference the theorem.
\response{\todo}

\item (p.~34, line~9) ``One source of matching invariants comes from the free magma'':
remind the reader that free magmas were defined in \S2.
\response{\todo}

\item (p.~38, line~18) ``Clearly the term being rewritten is in \dots'' should be
``Clearly, if the term being rewritten is in \dots''.
\response{\todo}

\end{requests}

\end{document}
